% Publication typesetting template. See COPYRIGHT_AND_LICENSE.md.
\documentclass[11pt,a4paper]{article}
\usepackage[margin=26mm,headheight=15pt,headsep=8mm]{geometry}
\usepackage{fontspec}
\setmainfont{Liberation Serif}
\setsansfont{Liberation Sans}
\setmonofont[Scale=0.87]{DejaVu Sans Mono}
\usepackage{amsmath,amssymb}
\usepackage{microtype}
\usepackage{setspace}
\setstretch{1.12}
\usepackage{longtable,booktabs,array,calc}
\usepackage{needspace}
\usepackage{graphicx}
\usepackage{fancyhdr}
\usepackage[authoryear,round]{natbib}
\usepackage{xcolor}
\definecolor{linkblue}{HTML}{244A66}
\PassOptionsToPackage{hyphens}{url}
\usepackage[unicode,colorlinks=true,allcolors=linkblue,bookmarksnumbered=true]{hyperref}
\usepackage{xurl}
\urlstyle{same}
\hypersetup{pdftitle={What Prediction Leaves Out: Relational Objecthood
and the Grounds of Purpose},pdfauthor={Degui
Qian},pdfsubject={A reconstruction of UEOT; version 1.0}}
\setlength{\parindent}{1.15em}
\setlength{\parskip}{2pt}
\setlength{\emergencystretch}{2em}
\setcounter{secnumdepth}{2}
\setcounter{tocdepth}{2}
\widowpenalty=10000
\clubpenalty=10000
\displaywidowpenalty=10000
\providecommand{\tightlist}{\setlength{\itemsep}{0pt}\setlength{\parskip}{0pt}}
\providecommand{\passthrough}[1]{#1}
\pagestyle{fancy}
\fancyhf{}
\fancyhead[L]{\footnotesize Qian: Objecthood and Purpose}
\fancyhead[R]{\footnotesize Research article \textperiodcentered{} 1.0}
\fancyfoot[C]{\small\thepage}
\renewcommand{\headrulewidth}{0.25pt}
\fancypagestyle{plain}{\fancyhf{}\fancyfoot[C]{\small\thepage}\renewcommand{\headrulewidth}{0pt}}
\begin{document}
\thispagestyle{plain}
\begin{center}
{\fontsize{21}{25}\selectfont\bfseries What Prediction Leaves Out:
Relational Objecthood and the Grounds of Purpose\par}
\vspace{11pt}
{\large Degui Qian\par}
\vspace{3pt}
{\small Independent Researcher\par}
\vspace{6pt}
{\small 11 November 2026 \textperiodcentered{} Version 1.0\par}
\end{center}
\begin{abstract}
\noindent When does a successful representation of a system justify
claims about what persists and what its activity is for? This paper
reconstructs Universal Evolutionary Object Theory around that question.
It defends a relational realism about organized processes and a
conditional account of purposive explanation, while withdrawing the
original programme's universal loop requirement and unrestricted
optimization claims. Its central argument concerns preservation:
transferring evidence from prediction to persistence or purpose requires
retaining the distinctions needed for the intended attribution, at its
stated level of precision. A finite regenerative model makes this
requirement concrete. The same four initial configurations admit two,
three, or four predictive classes as the recorded task changes; states
indistinguishable by their complete birth records can differ in repair
expenditure. In the same source, population growth can favour an action
that reduces the original individual's survival, and an optimal policy
need not be an internally executed programme. These results neither
establish an ontology nor naturalize purpose by themselves. They
identify where causal implementation, historical continuity, feasible
control, and independently supported evaluation must enter. Comparison
with computational mechanics, process philosophy, and organizational
teleology locates the proposed contribution in their constrained
integration, rather than in priority for their constituent ideas. The
resulting unity is a defeasible discipline of explanation, with explicit
conditions under which apparently compatible accounts fail to explain
the same organized system.
\end{abstract}
\noindent{\small\textbf{Keywords:} objecthood; predictive equivalence;
persistence; teleology; organizational realism; UEOT}
\vspace{10pt}

\section{Introduction: two questions about one
system}\label{introduction-two-questions-about-one-system}

Consider a device that loses a working instruction and restores it from
surviving internal copies. Compare it with a device whose working
instruction was never damaged. Suppose the two subsequently produce
exactly the same distribution of observable outputs. For predicting
those outputs, remembering the damage may be unnecessary. For
reconstructing this device's actual repair expenditure, the distinction
is indispensable. Now suppose an operating mode increases the expected
number of descendants but slightly reduces the original device's chance
of remaining intact. Describing that mode as promoting the system's
persistence is ambiguous. It matters whether the system means this
individual, the organization reproduced in its successors, or a
population counted at a later date.

These are questions about the entitlement to transfer an explanation. An
adequate account of outputs can leave out a real difference in the
conditions of their production. An adequate account of reproduction can
leave out the death of the reproducing individual. Neither account
becomes false because it is selective. Trouble begins when its success
is taken to establish a stronger conclusion about the object or its
purpose. The central claim of this paper is that object identification
and purposive attribution are distinct but mutually constraining
explanatory tasks. A representation adequate for one supports the other
only if, with any justified supplementary evidence, it preserves what
the intended attribution requires. Exact individual histories,
probability laws, and aggregate expectations can require different
information.

Universal Evolutionary Object Theory, or UEOT, originally approached
these questions through a much stronger programme. It proposed that
every persistent object depends on an internal self-maintaining loop,
that change reflects physical and informational drives, and that a
globally optimal direction and attractor could provide a naturalistic
basis for teleology. The original paper also described itself as a
research programme and contained qualifications about innovation and
changing environments. Those qualifications did not remove its explicit
universal assertions \citep{qianOriginal}. The present article retains
the question of how objects and purposes belong in one naturalistic
account while revising the proposed answer.

The revision is substantive. Persistence does not generally entail a
topological cycle. Predictive equivalence does not establish numerical
identity. An optimal-control calculation does not show that an organism
executes the calculated policy. Growth of a population is not survival
of an individual. Nor does successful maintenance, however accurately
described, establish an obligation that the maintained organization
ought to continue. The current UEOT materials acknowledge these
separations and provide conditional positive results, explicit
counterexamples, and a finite regenerative source in which several
explanatory relations can be checked together
\citep{qianPhilosophy2026, qianMathematics2026}. This article develops a
philosophical argument from those distinctions. It is not a summary of
the books or a claim that a proof assistant has established a
metaphysics.

The two positions defended here require different arguments. The first
is a selective relational realism: some organizations are real because
their relations support stable causal explanations across appropriately
varied encounters, even though no single predictive partition captures
everything about them. The realist commitment is defeasible and extends
beyond the mathematics of prediction. The second is a conditional
account of purposive explanation: a claim that activity is for an
outcome requires evidence connecting a relevant evaluative distinction
with the system's implementation or history. Such evidence can concern
regulatory control, a selected effect, or a designed function. These
grounds are not interchangeable, and they do not require a unique
numerical utility.

Their connection is the paper's principal contribution. A purpose
belongs to a specified bearer or organized activity. Changing the bearer
from an individual to a lineage may change the apparent purpose without
changing the underlying events. Conversely, attributing a function can
reveal distinctions that an otherwise adequate predictive representation
must retain. The question is therefore not simply which ontology is
predictive or which objective is optimal. It is whether the
representation of the bearer and the explanation of its activity can be
satisfied by the same process under the same restrictions. Separate
examples of a repairing system and an optimizing system do not establish
that one admissible system does both; a common realization is an
additional achievement.

The ingredients have substantial precedents. Computational mechanics
develops predictive equivalence and minimal representations. Process
philosophy understands persistence through organized change.
Organizational accounts of function and agency already connect
individuality, regulation, and normativity
\citep{shaliziCrutchfield2001, dupreNicholson2018, mossio2009, barandiaran2009}.
Recent work explicitly discusses regulatory teleology and the
persistence of the organized whole \citep{bich2024, cuciniello2026}. The
claim here is consequently narrower than a new foundation for all these
positions. It is an argument about their compatibility, sharpened by
paired failures and a restricted positive construction in one
mathematical source. Its value depends on whether that construction
exposes explanatory errors that would otherwise remain easy to overlook.

The argument proceeds from identification to persistence and then to
purpose. Section 2 explains which parts of the original programme
require revision. Section 3 distinguishes a predictive type from an
enduring bearer and develops the preservation condition. Section 4
examines organization, repair, and historical continuity. Section 5
shows why purposive explanation needs more than optimization, using
exact growth results from the same source. Section 6 explains the
limited unity achieved and its relation to the closest alternatives.
Section 7 addresses objections to the philosophical position and its
evidential reach.

\section{Reconstructing the original
programme}\label{reconstructing-the-original-programme}

\subsection{What survives revision}\label{what-survives-revision}

The enduring motivation of UEOT is that objects should not be treated as
inert bearers whose changes are philosophically secondary. In many
explanations, what persists is an organized capacity sustained through
material or informational replacement. The original programme asked
whether this perspective could also explain why some activities appear
directed toward maintaining or transforming their conditions of
existence. Those questions survive the failure of a universal answer. A
research programme can retain a problem while abandoning premises that
previously made its solution look immediate.

Three commitments are worth separating. First, relational organization
can be explanatorily important: what a component does depends on its
connections, available resources, and the larger process in which it
operates. Second, temporal maintenance is a phenomenon requiring
explanation, not a primitive guarantee supplied by naming an object.
Third, purposive language should be accountable to natural processes.
None entails that every object regenerates itself, that every persistent
system is an agent, or that nature possesses one privileged objective. A
stone, a repaired register, a reproducing lineage, and an autonomous
organism can share some explanatory features while differing in others.

The original distinction between Type I and Type II objects needs a fair
reading. It was presented as a contrast between state-dependent
organizational tendencies, not simply an immutable division of the world
into two kinds \citep[4-5]{qianOriginal}. Even so, a binary vocabulary
is too coarse to settle whether a particular process merely endures,
repairs damage, regulates a variable, reproduces its organization, or
acts with autonomous agency. The present account replaces reliance on
that contrast with explicit claims about these capacities. It does not
infer the strongest capacity from the weakest.

\subsection{What must be withdrawn or
narrowed}\label{what-must-be-withdrawn-or-narrowed}

The original loop postulate included nontrivial first homology of a
state space or interaction network. This is stronger than saying that a
persisting system has causal dependencies. A contracting flow, such as
\(\dot x=-x\) on a contractible state space, can remain in a bounded
region and approach a stable fixed point without nontrivial first
homology of that space. Stable state dynamics therefore do not entail
this state-space topology \citep[§§8-9]{qianMathematics2026}. The
example does not exclude a cyclic interaction network in some
realization of that flow; the original paper explicitly allowed that
alternative \citep[4]{qianOriginal}. A universal claim about an object's
network or regenerative organization needs a separate realization
argument. It cannot be inferred from state stability, and this flow
alone does not refute every such alternative. The reconstruction
withdraws the universal requirement because it lacks that general
justification, while retaining specific organizational hypotheses where
they are supported.

The physical and informational quantities denoted by \(\Pi\) and
\(\Phi\) also require separate measurement and interpretation. If only a
combined expression \(J=\Pi-\lambda\Phi\) is observed, the
transformation

\[
\Pi'=\Pi+\lambda\chi,\qquad \Phi'=\Phi+\chi
\]

leaves \(J\) unchanged. The decomposition is not identified by the
combined score alone. In the original discussion, informational
potential also shifted between uncertainty, adaptive novelty, and
evaluative significance. These uses cannot be unified by notation. The
current sources further supply a local obstruction to a universal smooth
two-dimensional representation. For the healthy care programme
introduced below, the public probabilities of retention, birth, and the
joint event of birth with damaged output are \((s,b,bq)\). Holding
\(\delta,\gamma\) fixed, their Jacobian with respect to \((s,b,q)\) has
determinant \(b>0\) on the legal open domain
\citep[§23.4]{qianMathematics2026}. These responses cannot all factor
there through one common smooth two-dimensional intermediate
representation. This is model-internal parameter variation, not a
demonstrated physical experiment, and it does not exclude every discrete
or set-theoretic encoding.

The language of a Global Optimal Direction and a Global Optimal
Attractor must likewise become conditional. An optimum is relative to an
objective, feasible actions, information, resources, and a time
criterion. Optimal actions can tie. An invariant distribution, a
conditional distribution among survivors, a deterministic limit set, and
a family of optimal-policy limits are different mathematical objects.
The current theory has results about particular members of these
categories; it does not establish a unique attractor toward which every
natural system progresses \citep{qianMathematics2026}. The original
paper's discussion of innovation already made an unchanging universal
optimum problematic, but its explicit convergence claims still require
withdrawal rather than reinterpretation as harmless metaphor.

Finally, the original extension from alignment with a global direction
to moral obligation does not survive the reconstruction. A mechanism's
conditions of operation can support judgments of functional success or
failure. They do not by themselves establish who ought to preserve it,
whether its continuation is desirable, or whose costs count. The new
article leaves moral justification to further argument. This is a change
in the scope of the theory, not merely a change of terminology
\citetext{\citealp[26-28]{qianOriginal}; \citealp{qianPhilosophy2026}}.

These revisions leave a positive task. If universal loops and objectives
cannot supply identity and purpose in advance, one must investigate the
relations that make a particular attribution defensible. The rest of the
paper argues that these relations constrain each other. A formal
distinction is useful here precisely because it can prevent an
attractive philosophical interpretation from acquiring more support than
its premises provide.

\section{Predictive identity and the reality of
organization}\label{predictive-identity-and-the-reality-of-organization}

\subsection{A representation answers a
question}\label{a-representation-answers-a-question}

Let a history \(h\) determine response laws \(K_i(h,\cdot)\) for
declared protocols \(i\). Its canonical predictive representation is the
family

\[
C(h)=\bigl(K_i(h,\cdot)\bigr)_i.
\]

Histories agree predictively when they yield the same relevant response
laws. Under the measurable and countability assumptions of the current
UEOT result, any statistic through which all these laws factor also
determines \(C\) almost everywhere with respect to the specified
reference measure. The conclusion is a precise minimality statement. It
is not the assertion that every conceivable conditional probability has
a pointwise canonical value, including on unobserved null histories
\citetext{\citealp[§§4.1-4.3]{qianMathematics2026}; \citealp{ueotSource2026}}.
The supplementary formal matrix supplies the exact source route as F01;
the common-source construction discussed next is F02.

This construction belongs to an established mathematical tradition.
Causal states identify histories by their distributions of futures;
transducer constructions extend predictive equivalence to input-output
processes \citetext{\citealp[causal-state definition and theorems
1-3]{shaliziCrutchfield2001}; \citealp[causal-equivalence and optimality
results]{barnettCrutchfield2015}}. Conditioning on supplied inputs does
not by itself establish causal intervention semantics. The current UEOT
source supplies its own declared protocol and noise restrictions, and
connects certain response laws to a common latent process. That
connection matters: a conditional future must arise from the process
that generated the past, rather than from an independently substituted
hidden state that changes the conditional joint law. Correct posterior
resampling is not excluded. But neither minimality nor source
consistency changes the kind of identity initially established. It is
equivalence for a specified response task.

The investigator chooses which question to ask. Once that question and
the model are fixed, its answer is constrained. A representation that
merges histories with different required future distributions is
inadequate, even if it is simpler or philosophically attractive.
Conversely, adding a historical distinction without predictive
consequences may make a statistic redundant for that task. These are
objective relations among specified models and questions. They do not
yet determine which questions reveal an enduring object.

Calling predictive classes types helps avoid a premature ontological
inference. Two numerically distinct devices can instantiate the same
predictive type. One device can move between predictive types after
damage. A coarser task can merge states that a richer task separates.
None of these facts implies that the device has ceased to exist, that
two devices have become one, or that objecthood is unreal. It implies
that identifying a response type and tracking a bearer are different
operations.

\subsection{One source, three legitimate
partitions}\label{one-source-three-legitimate-partitions}

The finite regenerative source makes the dependence on the task exact.
There are two internal programmes,
\(\theta_0=(\mathrm{care},\mathrm{care})\) and
\(\theta_1=(\mathrm{bloom},\mathrm{care})\), and two initial conditions,
healthy \(H\) and damaged \(D\). The ordered programme entries specify
actions after a healthy or damaged diagnosis. Bloom trades some
individual retention for increased birth probability; care is the
affordable executable action after damage and paid repair. Let \(B_t\)
count strict births across the whole actual population at census \(t\),
excluding every retention event. A newborn first executes at the
following census. The four initial configurations are \(\theta_0H\),
\(\theta_0D\), \(\theta_1H\), and \(\theta_1D\). Under Section 5's
probability conditions together with \(\gamma>\delta\) and \(s+b>1\),
the birth advantage exceeds the retention loss and care has expected
total output above one. In this original positive regime, the partitions
are those in Table 1 \citep[opening definitive argument, part
III]{qianMathematics2026}.

\Needspace{16\baselineskip}

\begin{longtable}[]{@{}
  >{\raggedright\arraybackslash}p{(\columnwidth - 2\tabcolsep) * \real{0.5000}}
  >{\raggedright\arraybackslash}p{(\columnwidth - 2\tabcolsep) * \real{0.5000}}@{}}
\toprule\noalign{}
\begin{minipage}[b]{\linewidth}\raggedright
Recorded task
\end{minipage} & \begin{minipage}[b]{\linewidth}\raggedright
Predictive classes of the four initial configurations
\end{minipage} \\
\midrule\noalign{}
\endhead
\bottomrule\noalign{}
\endlastfoot
First strict-birth count \(B_1\) &
\(\{\theta_0H,\theta_0D,\theta_1D\}\); \(\{\theta_1H\}\) \\
Joint birth records \((B_1,\ldots,B_T)\) for every finite \(T\) &
\(\{\theta_0H,\theta_0D\}\); \(\{\theta_1D\}\); \(\{\theta_1H\}\) \\
Those complete finite birth records together with first repair fee
\(R_1\) & Four singleton classes \\
\end{longtable}

Table 1. Task-dependent predictive distinctions within one
fixed-programme source. The result is a written proof with independent
exact verification in the source record; it is not represented here as a
complete dedicated Lean theorem. It concerns these fixed internal
programmes, not the entire class of external control policies.

The healthy and damaged starts of \(\theta_0\) can be coupled to have
the same entire birth record after the damaged body is restored. Their
initial repair expenditures differ: zero for \(H\) and three units for
\(D\). The complete birth process consequently omits a distinction
relevant to the explanation of resource use. At the same time, extending
the first-step birth question to all finite joint birth records
separates \(\theta_1D\) from \(\theta_0D\). Their initially identical
care actions lead to descendants that carry different programmes and can
act differently when healthy. The distinction appears in the second
census. At \((s,\delta,b,\gamma,q)=(3/4,1/8,3/8,1/4,1/3)\), the expected
second birth counts for \(\theta_0D\) and \(\theta_1D\) are \(27/64\)
and \(39/64\), respectively.

The third row requires attention to the whole task. Merely recording
\((B_1,R_1)\) produces three classes, since both damaged programmes
initially undergo the same payment and care action. Four classes result
when the repair fee is added to the complete finite birth-record task.
This detail is philosophically instructive. The record that resolves one
ambiguity need not resolve another; adequacy belongs to the joint
question actually posed.

All three partitions can be correct. The first does not falsify the
third, and the third does not reveal that the first was secretly
describing four predictive objects. They identify different equivalence
relations in the same source. Yet only the richer representation can
carry an explanation that explicitly distinguishes a paid recovery from
undamaged operation. If one uses the coarser birth representation to
argue that repair costs are constitutively irrelevant to the
organization, the inference fails. What was irrelevant to a selected
output task has been promoted into an unrestricted claim about the
process.

\subsection{The preservation
condition}\label{the-preservation-condition}

The logical point is elementary but consequential. Suppose \(q:S\to Z\)
is a representation of source situations and \(F\) is a feature required
by an explanation. A well-defined feature \(\bar F\) of the represented
states can satisfy \(F=\bar F\circ q\) only if \(F\) agrees whenever
\(q\) agrees. Conversely, in the simple set-theoretic case, agreement on
every fibre of \(q\) defines such a \(\bar F\) on its image. If the
explanation concerns histories, the same question must be asked of the
history representation, not just the instantaneous state map. For
measurable or probabilistic claims, the relevant regularity and
almost-everywhere qualifications must also be retained.

No new mathematical theorem is claimed by this observation. Its
immediate scope is an exact attribution of \(F\) from the stated
representation alone. A probabilistic or aggregate explanation need not
preserve each realized individual value. Suppose a quotient merges
repair fees zero and three, and independent evidence supplies equal
probability for the two cases. The quotient together with that evidence
can determine a fee distribution and expected fee \(3/2\), although it
cannot identify which fee this bearer paid. For a distribution,
expectation, or approximation, preservation concerns that stated target;
supplementary evidence belongs in the combined representation being
assessed.

The philosophical proposal is to use this distinction as a discipline on
explanatory transfer. Identify the bearer, explanatory contrast,
temporal scope, and basis of evaluation before claiming that a
convenient representation answers the question. If an exact persistence
claim requires identifying whether a bearer survived, the available
evidence must retain that distinction or supply a justified history map.
If an explanation identifies a particular paid restoration, subsequent
outputs alone can be insufficient. If optimal-action reasoning relies on
affordability, an observation class erasing the relevant difference
cannot support that same state-specific decision. These restrictions
allow useful aggregate explanations while blocking their unsupported
promotion to individual histories.

The condition is necessary, not an automatic warrant. Constant features
on a quotient can still be interpreted badly, the source model can be
false, and an evaluative distinction can be arbitrary. Preservation
prevents information already known to matter from being silently lost.
It does not manufacture the evidence that the information is true or
relevant. In particular, choosing the conclusion as a new feature would
make it formally preserved while doing nothing to justify it. The burden
is to establish the feature's role independently and then determine
whether the representation retains it.

There is a positive counterpart in the same source. A causally silent
physical tag can be erased while preserving the declared projections of
transactions, admissible observations, payments, and marked histories.
Raw bodies carrying different tags are not literally identical; the
relevant records and history laws agree under the source's projection.
Erasing the diagnosis can instead remove a distinction in affordability,
while preserving decoded instruction bits alone can lose the distinction
between an open and a closed execution gate \citep{ueotSource2026}. The
principle therefore does not demand maximal detail. It distinguishes
justified omission from omission that changes the target explanation.

\subsection{From representation to selective
realism}\label{from-representation-to-selective-realism}

Prediction alone does not deliver realism. The further argument is that
an organization deserves realist interpretation when its relations
explain stable differences across independently specified perturbations,
its realization can be traced, and competing representations can be
corrected by the same worldly process. The grounds are causal and
comparative. In the regenerative example, whether the gate is open and
whether the backups survive make differences to actual execution and
repair. These relations are not made true by preferring one vocabulary.
In a natural case, analogous claims would require observation and
intervention rather than inspection of stipulated source code.

This position is close to pattern realism without reducing all reality
to compression. Dennett's account already connects predictive patterns
and ontological commitment, while allowing that more than one useful
pattern can be found \citep{dennett1991}. The present argument concerns
the extra entitlement needed for a particular persistence or purpose
claim. It also respects the difference between structural knowledge and
a stronger ontological thesis emphasized in debates about structural
realism \citep{ladyman1998}. A mathematical quotient cannot settle that
debate by changing its name to an object.

Protocol relativity thus need not imply that organization is an
invention. Selecting the temperature to measure does not select its
value; selecting a response task does not select which histories are
equivalent. Nevertheless, a modeler can make a preferred classification
trivially correct by choosing a weak enough task. The realist inference
gains strength when the organization survives relevant changes in
measurement, when failed interventions are recorded, and when the model
risks being wrong. The result is pluralism under constraint, rather than
a universal algorithm for identifying the one true partition of nature.

\section{Persistence, repair, and the bearer of a
history}\label{persistence-repair-and-the-bearer-of-a-history}

\subsection{Organization is more specific than
recurrence}\label{organization-is-more-specific-than-recurrence}

Rejecting a universal topological loop leaves several distinct
explanatory possibilities. A dynamical system can remain within a domain
without repairing itself. A controller can return a variable to a range
without producing its own components. A network can regenerate
components while depending on external energy and trusted constraints.
These distinctions concern what actually sustains the capacity, not
merely the vocabulary used to describe repeated activity. A trajectory's
recurrence, a graph cycle, feedback in a controller, and closure among
constraints are therefore not interchangeable evidence.

The organizational approach offers a stronger biological hypothesis than
generic recurrence. Autopoiesis concerns a productive organization that
generates its components and realizes a unity; reproduction is not
simply substituted for that organization \citep{varela1974}. Closure of
constraints concerns interdependencies among relatively stable
constraints on material and energetic processes, with attention to their
timescales \citep[definitions 1-4 and §§5-6]{montevilMossio2015}. The
contracting-flow example in Section 2 neither instantiates nor refutes
those hypotheses. It separates stable state dynamics from nontrivial
state-space homology. An externally driven cycle can also repeat without
regenerating its driver, while a regulator can maintain a capacity
without following a closed geometric orbit. Thus trajectory topology,
network feedback, and regenerative dependence answer different
questions. Which organization a living system realizes needs its own
evidence.

The revised use of \(\Omega\) illustrates both the gain and the limit of
a more precise approach. Fix a persistence certificate, a horizon, a
family of permitted deletions, and a failure criterion. When failure is
monotone under deletion, its finite destructive sets can be
reconstructed from their minimal members. These minimal failures can
identify dependencies relevant to the certificate
\citetext{\citealp[§12.2]{qianMathematics2026}; \citealp[supplementary
matrix F05-F07]{ueotSource2026}}. The construction depends on those
choices and on monotonicity. Removing an inhibitor may activate a
backup, so an arbitrary deletion semantics need not be monotone. A
restricted intervention family can also have no destructive member at
all.

This causal reading of \(\Omega\) cannot quietly inherit the original
universal claim. If the permitted interventions remove only an
irrelevant input, their failure family says little about the robustness
of the remaining organization. Conversely, losing access to a readout
can prevent certification while leaving the process intact. A useful
dependency analysis must distinguish loss of evidence from loss of the
mechanism itself. The revised concept is valuable when its intervention
and certification conditions are informative; it is empty when those
conditions have been chosen merely to secure a desired verdict.

\subsection{A repair that is implemented rather than
stipulated}\label{a-repair-that-is-implemented-rather-than-stipulated}

The regenerative source supplies a modest positive case. Each body has
two instruction rows, with one active bit and two backup bits per row, a
main execution gate, and a bootstrap capability. The specified damage
flips the active bits and closes the gate while retaining the backups.
The repair routine reads the available copies, restores the active bits
by majority, and reopens the gate. Subsequent execution reads the
restored physical cells in the model. It does not consult an externally
supplied label declaring the correct programme \citep{ueotSource2026}.

Resources are part of the same event. A body begins with one unit and
receives five additional units per census. The specified damaged body
spends three on repair, one on care, and reserves two for possible
outputs. The model accounts for the complete six-unit budget. Outputs
inherit the restored code before undergoing the specified later damage.
The claim is consequently stronger than an interface that simply assumes
successful repair, available funds, and faithful inheritance
independently. Within this source, those relations are tied to the same
body and transaction.

It remains a finite mathematical model. Its physical cells and payments
are explicit model variables, not observations of a constructed device.
The bootstrap and writing operations are trusted primitives. The damage
leaves information in the backups. The external resource inflow is
prescribed. Later noise can still eliminate the body. These
qualifications determine what the example shows: a logically and
causally specified conjunction of repair, execution, and reproduction is
possible. It does not show that every organism has this architecture or
that an arbitrary damaged system can recreate information it has
entirely lost.

The accompanying failures make the positive claim informative. Cutting
the bootstrap prevents the specified repair. Insufficient funds at the
repair stage prevent the required paid operation. Erasing all programme
information cannot leave a single decoder that recovers two distinct
original programmes from the same remaining input
\citep{ueotSource2026}. These are failures of different enabling
conditions. None is repaired by observing that some abstract controller
could output the desired bit pattern. The problem is whether this body,
with these accessible resources and operations, can do so.

Here the preservation condition changes the explanation. In the named
raw-execution witness, two bodies have the same decoded instruction: the
open gate permits the specified bloom action, while the closed gate
blocks raw execution. This comparison fixes the diagnostic argument; it
is distinct from computing a new diagnosis or executing a complete
repair transaction. Code identity alone therefore does not determine
executable control. Two post-repair bodies can have identical birth
records while differing in the paid history that made those records
possible. Output identity alone does not recover that repair history. An
explanation of this recovery has to retain the causal route, even when
prediction of later outputs does not.

\subsection{Continuity, copying, and
branching}\label{continuity-copying-and-branching}

A repair story also contains a claim about whose activity continued. The
source constructs retention and strict birth as different events,
assigns their corresponding identities, and records parentage. The
labels track that event convention; names alone do not establish causal
continuity. A retained body and a newborn carrying the same programme
are not the same token merely because they behave alike. The historical
record permits questions about continued identity, lineage, and
reproducibility to have different answers. This is a mathematical
implementation of a distinction; it does not prove that its chosen
identity convention is the only legitimate one for natural organisms.

Process philosophy already treats material replacement and organized
continuity as central to biological persistence
\citep{dupreNicholson2018}. The present point is that a representation
can preserve a process's output organization while erasing the history
required by a particular continuity claim. Restoring the same
organization from a saved description might establish reproducibility.
Showing that a damaged bearer carried out the restoring operation might
establish continuity under an appropriate tracking relation. Neither
follows from matching future distributions. A causal-historical bridge
needs to specify which relations connect the earlier and later bearer
and why those relations are relevant to the inquiry.

Branching exposes the limit of a purely causal criterion. Suppose one
organization divides into two equally continuous successors. Both can
inherit its structure, memory, and activity. If both were numerically
identical to the predecessor, the symmetry and transitivity of identity
would make them identical to one another, contrary to their stipulated
distinctness. A lineage relation can branch; numerical identity cannot
do that same work without revision. The alternatives include selecting a
privileged continuation, treating the original as ending in division, or
acknowledging that the inquiry uses a different persistence relation.
Nothing in predictive equivalence selects among them.

This is especially important for teleology. A reproductive activity can
preserve an organization through new tokens while ending the individual
that performed it. A purpose explanation directed at lineage
continuation can be true without being an explanation of that
individual's continued existence. Recent discussion of organizational
teleology rightly treats diachronic identity as part of the
naturalization problem \citep{cuciniello2026}. The contribution here is
to make the relevant loss of identity information visible in the same
source used to discuss growth, rather than to claim that causation alone
solves every identity puzzle.

The resulting realism is limited but substantive. It permits recognition
of organized processes whose continuity is supported by causal history,
while allowing that some cases remain indeterminate or depend on a
stated tracking convention. It rejects an inference from behavioral
resemblance to numerical identity and an inference from causal descent
to uninterrupted individual survival. Those rejections become decisive
when a theory attempts to explain maintenance or purpose using a
representation that has discarded the bearer of the relevant history.

\section{The grounds and limits of purposive
explanation}\label{the-grounds-and-limits-of-purposive-explanation}

\subsection{What an objective does not
explain}\label{what-an-objective-does-not-explain}

An observed action sequence can often be made optimal by assigning it a
suitable reward. That possibility does not explain why the system acts
as it does or what its activity is for. The inverse reinforcement
learning literature already identifies underdetermination of reward by
optimal behavior \citep[§3.1]{ngRussell2000}. The current UEOT purpose
counterexamples make a related, restricted point: bare evidence and
feasible dynamics do not select a unique externally attached evaluation.
One formal construction simply pairs the same evidence value with
different objective functions. It must not be inflated into a proof that
two completely identical physical mechanisms have incompatible intrinsic
purposes \citetext{\citealp[§18.3,
N-TEL-01]{qianMathematics2026}; \citealp[supplementary matrix
F21]{ueotSource2026}}.

The distinction is between evidence for an evaluation and the mere
availability of a mathematical score. A regulatory target can have
support from the sensor-controller relation and the system's responses
to perturbations. A selected function can have support from the history
explaining a trait's retention. A designed function can have support
from a construction history. A contribution to a capacity can be
established by causal analysis of the organized system. These are
different explanatory achievements, as the contrast between etiological
and causal-role accounts of function makes clear
\citetext{\citealp[156-159]{wright1973}; \citealp[§III.3]{cummins1975}}.
A causal contribution alone need not be purposive; a further design,
historical, or regulatory ground has to support that attribution.
Calling all these achievements optimality suppresses the evidence needed
to distinguish them.

Independence here is evidential, not spatial or causal. A controller's
target need not come from outside the system, and a living organization
need not contain an explicit utility register. The requirement is that
the evidence for the relevant distinction does not consist only in
fitting a favourable score to the behavior to be explained. A directly
investigated repair mechanism may support the claim that its activity
restores a working organization even when several numerical rewards
rationalize that activity. Conversely, an externally chosen population
score may be precisely calculable without being represented as a goal by
any individual.

The finite source contains different kinds of relation. Its majority
routine restores damaged instruction cells under specified
perturbations. This is an implemented capacity. Calling restoration its
function has an additional ground in this designed model: the stated
construction task is to recover the pre-damage programme, whose
information is retained in the copies, and the routine is organized for
that task. This standard is not inferred from whichever output happens
to occur. It licenses a designed-function description, not an intrinsic
biological norm. The growth criterion instead evaluates population
consequences over time. Attainment does not show that the bearer
anticipates growth, values it, or regulates an estimate of it. The model
establishes implemented recovery and demographic performance; it does
not make the latter an intrinsic intention.

\subsection{Growth from the same
transactions}\label{growth-from-the-same-transactions}

Let \(s\) be the retention probability under care, \(b\) the birth
probability under care, \(\delta\) the retention cost of bloom,
\(\gamma\) its birth advantage, and \(q\) the probability that an output
is damaged. In the currently classified domain,

\[
0<\delta<s<1,\qquad b>0,\quad\gamma>0,\quad b+\gamma<1,
\qquad 0<q<1.
\]

The expected total outputs of the two executable actions are

\[
k_c=s+b,\qquad k_b=s-\delta+b+\gamma.
\]

The hereditary programmes introduced in Section 3 have exponential
factors

\[
\lambda_0=k_c,\qquad \lambda_1=(1-q)k_b+qk_c.
\]

These expressions come from the survival and birth transactions, the
output damage distribution, and the programme executed by each
diagnostic type. They are not rewards inferred from the observation that
a chosen programme wins. The distinction between retention and strict
birth remains in the actual histories even when the total population
count adds them \citep{qianMathematics2026, ueotGrowth2026}.

Under the source's full permitted public-history policy class and
canonical finite nonempty initial conditions, the least positive common
exponential upper factor for expected population is

\[
\rho_* = \max\{1,\lambda_0,\lambda_1\}.
\]

The exact result is expressed by bounds over every finite census. For
each fixed parameter choice and fixed nonempty finite canonical initial
population, there is a constant \(C\), independent of policy and census
time, such that \(\mathbb E_\pi N_t\le C\rho_*^t\) for every permitted
policy \(\pi\) and every finite \(t\). One specified stationary policy
has positive upper and lower constant multiples of the same exponential.
No smaller positive base supplies a common bound of that form. This does
not assume that every competitor has a growth limit. The corresponding
positive-root and limsup statements are analytical consequences of these
finite bounds, not separate new formal limit endpoints.

The two internal programmes attain precisely under the following
conditions:

\[
\begin{aligned}
\theta_0\text{ attains }\rho_* &\iff \gamma\le\delta\ \text{and}\ k_c\ge1,\\
\theta_1\text{ attains }\rho_* &\iff \gamma\ge\delta\ \text{and}\ \lambda_1\ge1.
\end{aligned}
\]

Each result is connected to that programme's own damaged founder, its
paid restoration, and its inherited execution. This is the positive
achievement of using one source. Repair is not demonstrated in one model
and optimality in an unrelated model with more convenient resources.
Nevertheless, the same-source connection does not erase the evaluation
choice: expected population is still the specified quantity being
compared.

The term \(1\) in \(\rho_*\) reveals why implementation cannot be
treated as a harmless technical detail. A damaged body's bloom request
would cost more than its available six units once repair and output
reserve are included. The actual transaction rejects that request,
preserves the body, returns the incoming resources, and does not expose
it to that transaction's noise. Repeating rejection therefore supplies a
population factor of one. This is a declared feature of this finite
source, not a natural law that inactivity prevents decay. If rejection
instead exposed the body to an independent mortality process, the action
rows and resulting optimization problem would change.

The permitted controller can use the public history and jointly
randomized requests; it cannot inspect the current world noise before
choosing or use hidden audit fields as observations. The comparison is
also supported by continuing resource inflows. It is not a theorem about
indefinite growth from one closed finite wallet. These premises identify
the optimization problem to which the conclusion belongs. Enlarging the
information set or replacing the rejected transaction with a different
physical event would change the problem rather than strengthen its
answer.

\subsection{Which persistence is
favoured?}\label{which-persistence-is-favoured}

On a healthy body, bloom decreases the original individual's retention
probability from \(s\) to \(s-\delta\) while increasing birth
probability from \(b\) to \(b+\gamma\). Where the gain is sufficiently
large, the population criterion favours bloom. The event can thus be
growth-promoting and individually more dangerous at the same time. A
statement that the action promotes persistence acquires a determinate
truth value only after its bearer and measure have been specified.
Population increase does not retrospectively make the dead original body
survive.

Even within the population evaluation, an exponential optimum differs
from the best count at a deadline. At the classified parameter point

\[
(s,\delta,b,\gamma,q)=(2/3,1/8,1/4,1/4,1/3),
\]

\(\lambda_1=\rho_*=1\). From the damaged founder, the internal
\(\theta_1\) programme has expected population \(11/12\) at every
positive census. A policy that repeatedly makes the rejected
damaged-body request retains one body on every path. Both have
exponential factor one, although their finite-time expectations differ
\citep{qianMathematics2026, ueotGrowth2026}. The comparison shows why an
asymptotic ranking cannot silently substitute for every practical
success criterion.

Nor does a higher expected factor imply that the preferred programme
survives in each finite realization. The source includes an exact
finite-population event in which the higher-growth programme disappears
while the lower-growth programme remains. Its selection theorem retains
the narrower original regime: the probability conditions above together
with \(\gamma>\delta\) and \(s+b>1\). For each specified share tolerance
\(0<\varepsilon\le1\) and error probability \(\eta>0\), it chooses a
finite positive horizon and then a finite number of independent cohorts.
Each cohort begins with one healthy founder of each programme and has
its own continuing resource account. At the chosen census, the
high-programme share is at least \(1-\varepsilon\) with probability at
least \(1-\eta\), counting extinction as failure. Independence concerns
the cohort noise laws, not funding alone. This is not almost-sure
selection over unlimited time in one fixed finite population
\citep[supplementary matrix F18-F19]{ueotSource2026}. Selection,
implemented action, and optimality therefore have separate evidential
burdens.

These examples support pluralism about purposive explanations, but not
unrestricted subjectivism. The survival probability, birth rate,
payment, and genealogy are fixed by the model. The evaluator cannot
reverse their values by changing a label. What remains open is which
relation grounds a particular purpose claim and which entity bears it.
Restoring the individual, transmitting a programme, increasing expected
descendants, and retaining a selected trait can all be naturalistically
describable without constituting one universal end.

\subsection{Functional standards and
reasons}\label{functional-standards-and-reasons}

Organizational accounts seek more than an observer's preference by
connecting functional contributions with the conditions that maintain
the organization. Mossio, Saborido, and Moreno give a substantial
account of such functions and their normative dimension
\citep{mossio2009}. Bich's regulatory development asks how activities
are coordinated in ways that matter to the organism's continued
organization \citep{bich2024}. These are genuine competitors to a view
that only externally assigned scores make teleology intelligible. The
present paper accepts that organization and regulation can supply
evidence unavailable from behavior fitting alone.

The force of the resulting standard remains contested. Corti's objection
concerns biological functional correctness: why does dependence of an
organization's existence on an activity establish how a component ought
to work, including a basis for malfunction? It is not merely a question
about moral goodness \citep[§3]{corti2023}. Causal effectiveness and
counterfactual dependence do not by themselves answer it. The present
model's repair standard is licensed by its specified design task and the
relation to the previously encoded programme. That warrants success or
failure relative to a designed function. It does not establish that this
standard is an intrinsic biological norm grounded only in the
organization's continued existence.

Four levels must consequently remain distinct: a descriptive causal
contribution; a functional attribution supported by design, formation
history, or an adequately defended regulatory account; intrinsic
biological normativity; and moral or practical reasons. The model
securely demonstrates the first and supports a designed-function
instance of the second. It does not resolve the third. Even if a
stronger biological account supplied that further ground, justification
of moral demands would remain an additional issue. This leaves the
strongest normativity objection open at its proper level instead of
pretending to answer it with a restriction on ethics.

The positive purposive claim is therefore specific: under the
independently identified design standard, the recovery routine is for
restoring executable code. The population calculation is a demographic
performance explanation unless further grounds make growth a function or
target of the bearer. Its precision does not itself supply those
grounds. Neither conclusion attributes consciousness. Autonomous agency
would also require further arguments concerning individuality,
interaction, and normativity of the kind already developed in the agency
literature \citep{barandiaran2009}. The distinction preserves a real
implemented accomplishment without turning every causally effective
process into a purposive one.

\section{What is unified, and what remains
plural}\label{what-is-unified-and-what-remains-plural}

\subsection{A shared constraint on explanatory
transfer}\label{a-shared-constraint-on-explanatory-transfer}

The two main arguments can now be joined. Predictive identification
relates histories through selected future records. Persistence relates
earlier and later bearers through an appropriate continuity claim.
Purposive explanation relates an activity to an outcome under a
function, control, or historical account. These relations need not have
the same extensions. A body can retain a predictive type after a costly
recovery; a newborn can share that type without being the earlier body;
a programme can increase descendants while diminishing its bearer's
retention. The unification proposed here consists in making these
differences constrain explanations of the same system.

Consider a constructed inference that tests the original UEOT ambition
to connect persistence, optimization, and purpose: the predictive
organization survives, the surviving organization optimizes growth,
therefore the object's repair is for its own optimal persistence. This
is a diagnostic reconstruction, not a quotation or an allegation that
computational mechanics or organizational teleology makes this
inference. Each clause may be defensible under some reading. The birth
quotient preserves the predicted organization; a policy attains the
growth factor; an individual undergoes repair. Yet the quotient may have
erased the repair expenditure, the optimal policy may be externally
specified rather than internally executed, and the counted population
may contain new identities. Persistence changes its bearer while
optimality changes its standard. The conclusion is not supplied by
placing the separately true premises beside each other.

The successful alternative has greater content than a warning. Let
\(R(x)\) mean that admissible system \(x\) performs the specified paid
recovery, and \(O(x)\) mean that it attains the stated benchmark.
Separate claims \(\exists x\,R(x)\) and \(\exists y\,O(y)\) do not imply
\(\exists z\,[R(z)\land O(z)]\). The source supplies the latter
conjunction for each named programme in its attainment region. The
common witness fixes the damaged founder, programme, backups, bootstrap,
transaction account, information permissions, and subsequent output
process. Repair and attainment are evaluated on that common realization.

The source remains constructed, and its resource inflow and rejection
rule remain modeling choices. Its added content is not that nature must
instantiate them. It is that the already specified six-unit account and
inherited execution impose potentially conflicting demands: damaged
bloom cannot be funded, repair alone does not ensure attainment, and an
external policy need not coincide with the stored programme. The result
resolves those demands under explicit parameter conditions without
changing the source for each desired property. This is a conditional
compatibility result; its empirical application still requires evidence
for the source assumptions.

This conjunction still warrants only the purpose claim supported by its
ground. It establishes restoration and growth attainment, not an
internal representation of growth. It establishes a causal relation
between programme and outcomes, not a moral reason to maximize them.
Thus the preservation condition joins the two arguments without
dissolving their differences. Object identification constrains whose
activity is explained; purpose attribution constrains which outcomes and
alternatives the representation must retain. Each can expose a
deficiency in a representation accepted for the other.

\subsection{Relation to the closest
alternatives}\label{relation-to-the-closest-alternatives}

The position continues a cybernetic insight that purpose can be
approached through organized relations among behavior, feedback, and
outcomes \citep{rosenblueth1943}. It also shares interventionism's
emphasis on the range over which explanatory relations remain stable
under changes \citep{woodward2000}. Its distinctive emphasis is the
compatibility of those relations with claims about the bearer's
identity. An input-output equivalence can be satisfactory until the
explanation asks which material or historical organization realizes the
response. A controller can be accurately described until its apparent
objective depends on an action it cannot perform.

Organizational teleology presents the closest challenge. Bich's account
concerns regulation that coordinates activities in relation to an
organism's organization, not just the presence of a mechanism
\citep[§§4-5]{bich2024}. Applying that criterion to the device is an
interpretation offered here, not a verdict attributed to Bich. Diagnosis
and the execution gate provide state-sensitive control, but this does
not demonstrate the wider biological organization or endogenous
normative ground his account seeks. External resources and trusted
components are not automatically disqualifying; their place in the
claimed organization still needs explanation. The preservation analysis
adds a specific test: a representation erasing the gate cannot preserve
this raw-execution distinction, even if it retains the decoded code. It
assesses a representation used by a regulatory explanation without
claiming to replace that explanation.

Cuciniello distinguishes the achievements and difficulties of minimal
and regulative organizational teleology, including the problem of how
regulatory goals and changing organizations are grounded
\citep[§§4-6]{cuciniello2026}. The present fixed-code recovery model
does not explain how novel goals arise or how an organization transforms
its own standards. Its separate retention and lineage records can
instead identify which bearer a proposed developmental or reproductive
explanation concerns. This is a narrower addition: an account that moves
from restored individual organization to transmission across generations
must justify that change of bearer. Neither competitor is accused of
neglecting mechanisms or history. Their stronger naturalization
questions remain open here; the proposed contribution is a way to test
whether a representation preserves the relations their explanations
invoke.

Other unifying programmes illustrate why scope matters. The free-energy
principle proposes a broad relation among organismic maintenance,
inference, and action; technical criticism has questioned particular
derivations and their assumptions \citep{friston2010, biehl2021}. This
article neither identifies its population objective with variational
free energy nor treats criticism of one formulation as a refutation of
every version. The comparable question is which bridge connects a
mathematical objective, a concrete generative process, and an
implemented policy. Naming the same functional across levels does not
itself establish that connection.

Likewise, results about causal emergence compare micro and macro
descriptions under a specified intervention distribution and information
measure \citep{hoel2013}. They do not settle that every coarser
description possesses greater causal power. They reinforce the need to
state what a comparison measures. The present view allows multiple
useful scales while requiring a representation to preserve the causal
differences on which the particular explanation depends.

This form of unity is methodological and philosophical. It does not
derive quantum dynamics, biological autonomy, and social institutions
from a shared equation. In each domain, the realization of an
organization, the legitimacy of interventions, and the ground of an
evaluative distinction require further inquiry. UEOT's current
mathematical closure is closure of declared conditional claims and their
boundaries \citep{ueotStatus2026}. Its current physics work retains
domain-specific assumptions and unresolved empirical responsibilities
\citep{ueotPhysicsStatus2026}. The paper's proposed unity can be
assessed without accepting a completed universal natural theory.

\section{Objections and consequences}\label{objections-and-consequences}

\subsection{The realist step remains
contestable}\label{the-realist-step-remains-contestable}

An empiricist can accept every mathematical result in this paper while
declining its realism. The finite source is a stipulated model; proving
relations within it does not show that nature contains the organizations
it depicts. Even if a natural system were accurately modeled, predictive
and intervention success might be treated as grounds for accepting the
model's observable consequences rather than its ontology. This objection
identifies a real limit. The argument does not deductively force realism
on an opponent who rejects explanatory inference to unobserved
organization.

The positive case is comparative. If a proposed organization explains
both ordinary behavior and changes following the loss of a gate, backup,
or resource, while a rival preserving only ordinary outputs does not,
there is a reason to treat the organization as more than a convenient
summary of that output record. The same relations account for
independently varying effects. The inference is defeasible because
another realization could explain them, the intervention may have
additional effects, or the purported continuity may be misidentified.
Realism here means accepting the explanatory organization to the extent
justified by those relations, with revision when their support fails.

This commitment is more limited than an inventory of intrinsically
bounded substances, but it is not merely a renaming of prediction. It
recognizes causal organization and historical dependence as possible
features of the world that a predictor can successfully ignore. A purely
predictive type can be an excellent instrument without being the
complete bearer of the explanation. The paper would lose its
philosophical point if it treated all representations as interchangeable
whenever their ordinary observations matched.

\subsection{Preservation can look
tautological}\label{preservation-can-look-tautological}

The preservation condition may appear to say only that an explanation
must retain what it needs. Taken alone, it does. Its value cannot lie in
the elementary fibre argument or in the introduction of another abstract
criterion. The substantive work is determining which distinctions a
seemingly successful explanation has actually erased, why they matter to
its conclusion, and whether retaining them is compatible with the other
desired claims.

The paired examples supply that work. Complete birth prediction does not
require the initial repair fee, although reconstructing this bearer's
paid-repair history does. Code equality does not require equality of
execution gates, although recovering the specified raw-execution
distinction does. Expected population combines retained and new
individuals, although a continuity explanation distinguishes them.
Growth attainment does not preserve finite-deadline rankings, although a
decision concerning a deadline requires them. These omissions are easy
to miss because the representation is successful in each of its original
tasks. The argument identifies a systematic route from legitimate
compression to illegitimate transfer.

The positive source construction adds a non-vacuity requirement. It is
insufficient to demand all desirable properties if they cannot coexist.
In the stated parameter regions, the paid repair, inherited internal
execution, and population attainment can coexist. Outside those regions,
failure of the two named programmes does not rule out every possible
repairing strategy; a permitted policy can, for example, pay for an
initial repair and then follow another rule. The result is an assessment
of particular connections, not a declaration that any bridge one wants
can be added. A reader may still judge this to be methodological
synthesis rather than a new metaphysics. That is a defensible
classification of the contribution, provided the concrete explanatory
discriminations are acknowledged.

\subsection{A model built to repair is a weak case for natural
purpose}\label{a-model-built-to-repair-is-a-weak-case-for-natural-purpose}

The strongest objection to the purposive reading is that a deliberately
constructed error-correcting model displays the investigator's design,
not the origins of natural purposiveness. The bootstrap, preserved
backups, resource inflow, and programme structure already encode many
conditions of success. A construction that starts with them cannot
explain how an organism acquires them or why its organization
constitutes a norm for itself. The present argument accepts this
limitation.

What the model establishes is more local: the presence of enabling
conditions does not make the implemented recovery fictitious.
Restoration actually changes damaged execution conditions, and the
source identifies failures when the necessary means are lost. A system
can perform some of its own maintenance while depending on supports it
does not produce. Treating every dependency as disqualifying
self-maintenance would make the relevant contrast disappear, since the
question concerns which conditions the organization itself restores.
Treating that limited restoration as complete autonomy would be the
opposite mistake.

A naturalized attribution therefore requires an additional account of
how the restoring organization is established and sustained in the
target. Selection history, developmental organization, or ongoing
regulation may provide different grounds. Their adequacy cannot be read
from a favourable expected-growth number. Nor must all purposive
explanations await one universal theory of intrinsic value: a
causal-role explanation can succeed on its stated terms while an account
of autonomous normativity remains disputed. The paper defends that
limited success and the need to identify which further claim is being
attempted.

\subsection{The framework risks conflating unlike
systems}\label{the-framework-risks-conflating-unlike-systems}

If every persistent organization is admitted, a rock, a thermostat, a
repaired register, and an organism may seem to become equivalent
examples of purpose. The objection would be decisive if the account
inferred agency from persistence. It does not. Enduring within a domain
establishes neither an error-sensitive response nor an internally
implemented repair. Regulation establishes neither production of the
regulating components nor a selection history. Reproduction establishes
neither conscious intention nor justified action. The relations can
overlap without forming a single ascending scale on which each system
has one score.

The finite device is consequently evidence for a restricted
compatibility claim, not a definition of life. A passive object can be
real without pursuing a goal. An organism can have reproductive
functions whose effects conflict with its immediate survival. An agent
can deliberately undertake a finite activity while accepting that it
will not preserve either itself or its organization indefinitely. The
vocabulary should track those differences rather than declare them
deficient approximations to one optimum.

The same restraint matters for human meaning. Understanding the material
and social conditions of a commitment can explain what carrying it out
requires and who bears its costs. It does not show that the commitment
is justified because it maintains a larger system. A reason to oppose a
destructive institution need not be invalidated by the institution's
capacity to reproduce itself. The reconstructed UEOT can illuminate the
feasibility and consequences of commitments while leaving their
justification to arguments about persons, harms, obligations, and
alternatives. That boundary preserves explanatory continuity without
deriving ethics from demographic success.

\subsection{The discriminating consequences are real but not yet
empirical
findings}\label{the-discriminating-consequences-are-real-but-not-yet-empirical-findings}

A philosophical framework should change how a claim can fail. This one
supplies several discriminations. If two models fit the same birth
record but disagree about repair expenditure, birth data alone cannot
choose between their resource explanations. If restoring output requires
an external overwrite after the internal bootstrap is removed, the
result does not confirm the original claim of internally implemented
recovery. If a policy succeeds only by inspecting current noise, it does
not establish success within the original information restriction. If a
tracking scheme counts a replacement as a retained token, its population
totals cannot independently verify uninterrupted individual persistence.

An empirical application could investigate these differences in a device
with inspectable memory and controlled faults. Its output task, internal
restoration operations, energy account, and identity tracking would
first need independent specification. One could then compare ordinary
output fit with the effects of disabling the repair path, modifying
accessible resources, or changing the allowed action interface. A
recovered output would be evidence for different claims depending on
whether it came from internal correction, external replacement, or a
previously hidden backup. This is a proposed research design
illustrating the argument, not an experiment performed for this article.

Biological application would be harder because interventions can damage
several capacities at once and boundaries change during development. The
difficulty does not license pretending that a successful abstract fit
resolves the realization question. It identifies why evidence for
regulation, organizational continuity, and function often has to be
assembled from several sources. A framework has empirical relevance when
it clarifies which observations would distinguish competing
explanations, even if its present contribution is to logical and
conceptual evaluation rather than a new measured effect.

The elementary factorization condition is analytic; its logical opposite
is not a useful empirical falsifier. What can be wrong is the judgment
that an explanation requires a certain distinction, that the chosen
representation actually preserves it, or that a natural process realizes
the proposed mechanism. The realist interpretation and the grounds of a
purpose attribution are also contestable. Beyond those failures, the
synthesis may add too little to established accounts to merit a
distinctive philosophical position. Formal precision helps expose these
questions but does not settle them. Their assessment requires comparison
with alternatives and criticism from readers outside the project.

\section{Conclusion}\label{conclusion}

The revised UEOT can sustain a philosophical position without a
universal topological loop, a unique optimum, or a derivation of
obligation from persistence. Organized processes can be real bearers of
causal and historical explanation even when their best predictive
descriptions depend on the task. Implemented restoration establishes a
causal capacity; an independently supported design, historical, or
regulatory account can additionally license an appropriate purpose
attribution. The finite example supplies a designed restoration function
and a separate demographic performance result, leaving intrinsic
biological normativity unresolved.

The two claims are connected by the distinctions an explanation must
preserve. In the finite source, predictive equivalence can erase paid
recovery, decoded information can erase executable control, and
population growth can erase the fate of the original individual. The
positive construction shows how some of these relations can instead be
satisfied together, under explicit assumptions, by the same damaged
founder and its successors. That conjunction establishes compatibility
and disciplined transfer, not a universal natural purpose.

The resulting unity is a commitment to explain the same organization
across descriptions while allowing different explanations to answer
different questions. Its philosophical merit lies in making the passage
between those questions answerable to causal realization, history, and
the grounds of evaluation. Its remaining ambition is to improve the
explanation of particular systems, and to let those explanations decide
how far the unity extends.

\section*{Sources, evidence, and
assistance}\label{sources-evidence-and-assistance}
\addcontentsline{toc}{section}{Sources, evidence, and assistance}

This article uses the historical UEOT paper and the maintained
philosophy and mathematics books at repository commit
\texttt{3d443ec29eb17af1fc00699a74e2d42d73516d81}, inspected on 11
November 2026. Mathematical references distinguish written arguments,
existing formal endpoints, and source-matched archived compilation
evidence. No new Lean theorem, complete Lean rebuild, physical
experiment, or empirical validation was performed for this article. The
complete task-dependent 2/3/4 classification and the local rank-three
response example are written mathematical results with exact source
checks, not newly supplied dedicated formal endpoints.

The accompanying formal claim-evidence matrix gives exact files,
endpoint names, assumptions, and exclusions. Its source routes are
F01-F02 for predictive minimality and the common conditional source;
F03-F07 for invariant limiting structures, homology, failure families,
and access; F08-F10 for execution, tag projection, and the birth
classifications; F11-F12 for paid repair and its failures; F13-F17 for
growth and internal attainment; F18-F20 for selection and marked
identity; F21 for purpose nonidentification; F22-F26 for tied
optimization, limiting structures, and the dual-drive boundaries; and
F27 for current scientific status. The general claim matrix separately
marks this article's philosophical interpretations. These supplements
belong with the article when its technical support is assessed.

Degui Qian is the author of the original UEOT programme and the named
author of this manuscript. ChatGPT assisted with research, drafting,
literature checking, and technical preparation. One coordinating
assistant wrote the manuscript; three native research agents separately
audited historical and philosophical claims, mathematical scope, and
prior art, and reviewed the complete draft. These are internal
AI-assisted reviews, not independent human peer review. No external
reviewer endorsement is asserted. The author remains responsible for the
scholarly claims and any submission decision.

Copyright © 2026 Degui Qian and respective rights holders, to the extent
applicable. This publication's original research prose is offered under
the Creative Commons Attribution 4.0 International licence. Previously
licensed source expression retains its prior terms. Third-party articles
are cited and paraphrased; their texts and figures are not relicensed or
reproduced here. The publication notice and third-party inventory
accompany the source package.

\clearpage
\begingroup
\setstretch{1.0}
\small
\setlength{\bibsep}{6pt}
\bibliographystyle{plainnat}
\bibliography{references}
\endgroup
\end{document}
